% GENERATED FILE. Edit canonical lessons, then run npm run generate:books.
% canonical-source-hash: fnv1a64:798d034a49bccf2b
% canonical-lessons: SA-C217-raktah, SA-C217-arunah, SA-C217-kapilah, SA-C217-samah, SA-C217-visamah

\chapter{Red, Reddish, Tawny, Equal, Uneven}
\label{ch:sa-red-reddish-tawny-equal-uneven}
\hlchaptermodality{217}

\begin{chapteropening}
\textbf{By the end of this chapter:} \emph{I can say red, reddish, tawny, equal and uneven.}

Complete the last lesson of chapter 217: I can say red, reddish, tawny, equal and uneven.
\end{chapteropening}

\section[\texorpdfstring{\sk{रक्तः} (raktaḥ)}{रक्तः (raktaḥ)}]{\sk{रक्तः} (raktaḥ) — red}

\label{lesson:SA-C217-raktah}



\begin{quote}
Before the new one: say the Sanskrit for wise, then the Sanskrit for intelligent.
\end{quote}

\subsection*{You'll want to know: \sk{रक्तः}}
\textbf{\sk{रक्तः}} — \emph{raktaḥ} — \textquotedblleft{}red\textquotedblright{}.

\begin{grammarlens}[title={What you've built}]
One more word for this chapter.
\end{grammarlens}

\subsection*{Your turn}
\begin{itemize}
\raggedright
  \item \emph{Say it:} \emph{raktaḥ}
  \item \emph{Say it:} \emph{raktaḥ}, once more
  \item \emph{Say it:} say \emph{raktaḥ}
  \item \emph{Recall:} say the Sanskrit for wise, then the Sanskrit for intelligent, then say \emph{raktaḥ} again
\end{itemize}

\begin{tcolorbox}[breakable,colback=teal!4,colframe=teal!35!black,title={Before you move on}]
What does \sk{रक्तः} mean? (\textquotedblleft{}Red\textquotedblright{}.) Say it once more.
\end{tcolorbox}
\section[\texorpdfstring{\sk{अरुणः} (aruṇaḥ)}{अरुणः (aruṇaḥ)}]{\sk{अरुणः} (aruṇaḥ) — reddish, dawn-red}

\label{lesson:SA-C217-arunah}



\begin{quote}
Before the new one: say the Sanskrit for intelligent, then the Sanskrit for red.
\end{quote}

\subsection*{You'll want to know: \sk{अरुणः}}
\textbf{\sk{अरुणः}} — \emph{aruṇaḥ} — \textquotedblleft{}reddish, dawn-red\textquotedblright{}.

\begin{grammarlens}[title={What you've built}]
One more word for this chapter.
\end{grammarlens}

\subsection*{Your turn}
\begin{itemize}
\raggedright
  \item \emph{Say it:} \emph{aruṇaḥ}
  \item \emph{Say it:} \emph{aruṇaḥ}, once more
  \item \emph{Say it:} say \emph{aruṇaḥ}
  \item \emph{Recall:} say the Sanskrit for intelligent, then the Sanskrit for red, then say \emph{aruṇaḥ} again
\end{itemize}

\begin{tcolorbox}[breakable,colback=teal!4,colframe=teal!35!black,title={Before you move on}]
What does \sk{अरुणः} mean? (\textquotedblleft{}Reddish, dawn-red\textquotedblright{}.) Say it once more.
\end{tcolorbox}
\section[\texorpdfstring{\sk{कपिलः} (kapilaḥ)}{कपिलः (kapilaḥ)}]{\sk{कपिलः} (kapilaḥ) — tawny, brown}

\label{lesson:SA-C217-kapilah}



\begin{quote}
Before the new one: say the Sanskrit for red, then the Sanskrit for reddish.
\end{quote}

\subsection*{You'll want to know: \sk{कपिलः}}
\textbf{\sk{कपिलः}} — \emph{kapilaḥ} — \textquotedblleft{}tawny, brown\textquotedblright{}.

\begin{grammarlens}[title={What you've built}]
One more word for this chapter.
\end{grammarlens}

\subsection*{Your turn}
\begin{itemize}
\raggedright
  \item \emph{Say it:} \emph{kapilaḥ}
  \item \emph{Say it:} \emph{kapilaḥ}, once more
  \item \emph{Say it:} say \emph{kapilaḥ}
  \item \emph{Recall:} say the Sanskrit for red, then the Sanskrit for reddish, then say \emph{kapilaḥ} again
\end{itemize}

\begin{tcolorbox}[breakable,colback=teal!4,colframe=teal!35!black,title={Before you move on}]
What does \sk{कपिलः} mean? (\textquotedblleft{}Tawny, brown\textquotedblright{}.) Say it once more.
\end{tcolorbox}
\section[\texorpdfstring{\sk{समः} (samaḥ)}{समः (samaḥ)}]{\sk{समः} (samaḥ) — equal, even}

\label{lesson:SA-C217-samah}



\begin{quote}
Before the new one: say the Sanskrit for reddish, then the Sanskrit for tawny.
\end{quote}

\subsection*{You'll want to know: \sk{समः}}
\textbf{\sk{समः}} — \emph{samaḥ} — \textquotedblleft{}equal, even\textquotedblright{}.

\begin{grammarlens}[title={What you've built}]
One more word for this chapter.
\end{grammarlens}

\subsection*{Your turn}
\begin{itemize}
\raggedright
  \item \emph{Say it:} \emph{samaḥ}
  \item \emph{Say it:} \emph{samaḥ}, once more
  \item \emph{Say it:} say \emph{samaḥ}
  \item \emph{Recall:} say the Sanskrit for reddish, then the Sanskrit for tawny, then say \emph{samaḥ} again
\end{itemize}

\begin{tcolorbox}[breakable,colback=teal!4,colframe=teal!35!black,title={Before you move on}]
What does \sk{समः} mean? (\textquotedblleft{}Equal, even\textquotedblright{}.) Say it once more.
\end{tcolorbox}
\section[\texorpdfstring{\sk{विषमः} (viṣamaḥ)}{विषमः (viṣamaḥ)}]{\sk{विषमः} (viṣamaḥ) — uneven, odd}

\label{lesson:SA-C217-visamah}



\begin{quote}
Before the new one: say the Sanskrit for tawny, then the Sanskrit for equal.
\end{quote}

\subsection*{You'll want to know: \sk{विषमः}}
\textbf{\sk{विषमः}} — \emph{viṣamaḥ} — \textquotedblleft{}uneven, odd\textquotedblright{}.

\begin{grammarlens}[title={What you've built}]
One more word for this chapter.
\end{grammarlens}

\subsection*{Your turn}
\begin{itemize}
\raggedright
  \item \emph{Say it:} \emph{viṣamaḥ}
  \item \emph{Say it:} \emph{viṣamaḥ}, once more
  \item \emph{Say it:} say \emph{viṣamaḥ}
  \item \emph{Recall:} say the Sanskrit for tawny, then the Sanskrit for equal, then say \emph{viṣamaḥ} again
\end{itemize}

\begin{tcolorbox}[breakable,colback=teal!4,colframe=teal!35!black,title={Before you move on}]
What does \sk{विषमः} mean? (\textquotedblleft{}Uneven, odd\textquotedblright{}.) Say it once more.
\end{tcolorbox}
